% Rebuttal and closing pages, revised from a first-time reader's perspective.
\Page{12 / 评审回应}{收到评审后，先处理最关键的问题}{找到已有答案，只为缺口补研究，再组织回应。}
\Lead{作者首先要判断：哪些疑虑能用现有证明和实验回答，哪些确实需要补做研究。}
\Fig{\resizebox{.95\textwidth}{!}{\RebuttalFlow}}{从论文和评审出发，合并相同问题，找到对应证据。需要补研究时先补齐，再起草、检查和修改回应。}

\Sub{先找答案，再决定补什么}
同一个问题可能被几位审稿人反复提出。EAR 先把这些意见放在一起，找到论文中的对应论据，再判断缺的是解释、证明还是实验。已有材料能够回答的，就组织清楚；确实缺证据的，才安排补充研究。同一项补充实验可以用于回答多条相关意见。\src{1,5,9}

\Sub{把时间花在最影响评价的问题上}
社区公开的回应方法和作者经验提供起点。系统比较不同的回答顺序、证据选择和写法，优先处理影响论文判断的疑虑。实现时，将这个任务写成一个优化目标：
\[
\max_{\theta}\ \mathbb E[\Delta R\mid\mathrm{Context},\theta]
\quad\text{subject to}\quad H\le H_B,\ C\le B,\ T\le D.
\]
这里的 Context 就是论文、评审和已有证据；$\theta$ 是回应策略，$\Delta R$ 是希望改善的评分，$D$ 是截止时间。含义很直接：在可用的人时、费用和期限内，选择最有希望解决关键疑虑的回应方式。

\begin{insight}
\textbf{模型给建议，作者检查关键内容}\quad 模型检查问题有没有回答、论据有没有依据，并预测评分可能上涨、不变或下降。这些反馈用于修改草稿；新增技术结论和最终回应由作者检查。
\end{insight}

\Page{13 / 回应策略优化}{保留好草稿，继续解决剩下的问题}{从社区已有方法开始，每轮都接着上一轮的结果改。}
\Fig{\resizebox{.95\textwidth}{!}{\RebuttalPopulation}}{尝试不同写法，比较模型评价，保留较好的草稿。下一轮带上尚未回答的问题和修改建议，继续完善。}

\Sub{先尝试几种有依据的写法}
初始策略来自公开的回应方法和作者经验：可以先解释已有证据，也可以先回应审稿人最关心的问题，或先安排缺失的验证。它们使用同一份论文和评审材料，区别在于如何组织工作和回答。

\Sub{每轮留下草稿、问题和建议}
系统保留当前较好的回应，并记下哪里还没讲清、哪些证据还不足、下一步应该怎么改。下一轮直接接着这些材料推进。实现上，这是黑盒优化：通过比较候选回应和读取反馈来调整策略，无需计算评分对写作策略的导数。

\Sub{记录中确实发生了策略调整}
一份匿名两轮记录中，第一轮选中的草稿未达到内部模型的评价要求。系统保留草稿和修改建议，第二轮换用另一策略后通过评价。这展示了系统会根据反馈调整做法，记录中的结果来自内部模型检查。

草稿通过内容评价后，系统停止尝试本轮其余策略，再检查回应是否忠实于论文和证据、表述是否准确。发现问题就继续修，检查完成后交给作者审阅。\src{1,10}

\Page{14 / E4：回应效果预测}{模型能预测回应后的评分变化吗？}{用已经发生的评审案例，检查模型的判断有没有依据。}
\Lead{模型读取论文、初始评审和作者当时的回应，预测评分会上涨、不变还是下降。}
最终评分只用来核对预测。初始评分和最终评分的差，给出真实变化标签；每个案例单独输入模型，避免不同案例之间混入信息。测试结果如下。\src{1,8,10}

\begin{center}\small
\begin{tabularx}{\textwidth}{P{32mm}Y P{21mm}}
\toprule
\textbf{历史案例} & \textbf{预测任务} & \textbf{Macro-F1}\\
\midrule
36 例均衡子集 & 区分提分与不变，各 18 例 & .805\\
57 例完整测试集 & 区分提分、不变和降分 & .503\\
同一 57 例来源集 & 来源报告的提分／不变二类指标 & .755\\
\bottomrule
\end{tabularx}
\end{center}
\captionof{table}{DeepSeek V4 Pro 的记录结果。Macro-F1 分别计算每个类别的 F1 后取平均，取值越接近 1 越好。三行对应不同测试设置。}

\Sub{能区分提分与不变，识别降分仍然较弱}
在 36 例均衡测试中，模型区分提分和不变的 macro-F1 为 .805。57 例完整测试集中包含 6 例降分情况，三类 macro-F1 为 .503，其中降分类的 F1 为 0。这提示系统不能只看一个总体分数，还要检查模型在哪一类判断上容易出错。

\Sub{把反馈用于改草稿，也检查它的依据}
EAR 将评分变化预测与证据检查一起用于比较回应。测试中，进一步修改提示没有得到更好的留出集成绩，因此保留原始零样本提示，没有修改模型权重。

\begin{insight}
\textbf{这些数字回答的是预测能力}\quad 测试使用已经写好的历史回应。系统生成新回应后，真实审稿人是否因此提高评分，需要用另一项实验检验。
\end{insight}

\Page{15 / 对研究者意味着什么}{让研究持续推进，把关键判断留给人}{从筛选问题到修订论文，再到组织回应，EAR 接手反复执行的工作。}
研究会反复经历同一类过程：查一轮文献，发现缺口；做一次尝试，遇到问题；补完证据，再修改结论。每次重新整理背景、安排下一步和检查遗漏，都会占用研究者的时间。EAR 把这些步骤接起来，让已经得到的材料和反馈能够继续用于下一步。

\Sub{找问题：让新证据改变投入方向}
搜索模块持续比较候选与已有工作。记录中，新找到的近似工作使一个候选被放弃，并留下退出理由。研究者可以先检查这些关键差异，再决定是否投入更大规模的推导和实验。

\Sub{做研究：把尝试、检查和修订接起来}
内循环保留证明、代码、失败路线和评阅意见，按具体问题继续修。固定的 28 个稿件中，当前版本通过模型评阅的数量从初轮 3 个增至第三轮修订后的 12 个。外循环再修改研究策略：在展示到 G3 的阶段中，当代选中策略四次尝试有三次通过；每个通过结果的研究调用，比同在 G3 运行的初始策略少 52.8\%。\src{10}

\Sub{写回应：把精力放在尚未解决的疑虑上}
回应模块合并重复问题，查找已有依据，必要时安排补充研究，再比较和修改草稿。每一轮保留有效内容和下一步建议，作者可以直接检查关键证据、技术结论和最终表述。

\begin{insight}
\textbf{一次判断，带动下一段工作}\quad 人给方向、设标准、检查关键结果；Agent 接着检索、尝试和修订。研究者回来时，有新的结果和清楚的修改记录，可以据此决定下一步。
\end{insight}

\noindent 本报告展示了流程、模型评阅和调用结果。实际节省多少主动人时，还需要在使用中直接记录。
